\documentclass[11pt]{article}
\usepackage{acl}
\usepackage{times}
\usepackage{latexsym}
\usepackage[T1]{fontenc}
\usepackage[utf8]{inputenc}
\usepackage{microtype}
\usepackage{inconsolata}
\usepackage{booktabs}
\usepackage{amsmath}
\usepackage{graphicx}
\usepackage{float}

\setlength\titlebox{6cm}

\title{Midway Report: Citation-Grounded Scientific QA with Local Retrieval-Augmented Generation}

\author{
  Jaswanth Pinnepu \\
  Boston University \\
  \texttt{jaswanth@bu.edu}
  \And
  Zukhriddin Fakhriddinov \\
  Boston University \\
  \texttt{zfmsai@bu.edu}
  \And
  Atharv Gulati \\
  Boston University \\
  \texttt{atharvg@bu.edu}
  \And
  Aayush Kumar \\
  Boston University \\
  \texttt{aayushks@bu.edu}
}

\begin{document}
\maketitle

\begin{abstract}
We study hallucination reduction in scientific question answering with a fully local retrieval-augmented generation (RAG) pipeline. Our setting uses QASPER as the benchmark, Phi-3.5-mini-instruct as the main generator, and shared Boston University SCC infrastructure for reproducible experiments. By the midway point, we have completed the data pipeline, chunking strategies, retrieval stack, and the proposal's core generation ablations A--E, along with three additional baselines: parametric-only generation, BM25-only retrieval, and random retrieval. On the original 100-question validation sample, the strongest core condition is section-aware chunking with hybrid retrieval and baseline prompting (Condition C), which reaches token F1 0.3947, BERTScore F1 0.8946, hallucination rate 0.39, and supported-answer rate 0.79. We subsequently extended the strongest settings to full validation and added an NLI-based citation check: on 1005 validation questions, the best overall Phi baseline (Condition D) reaches token F1 0.3238, while citation forcing (Condition E) reaches token F1 0.3215 with citation precision 0.3230, citation rate 0.9274, and NLI citation precision 0.1416. Citation forcing improves explicit attribution behavior, but it still reduces grounding relative to the strongest baseline setting.
\end{abstract}

\section{Introduction}

Our project asks how much hallucination in a small, locally hosted language model is caused by retrieval quality versus generation behavior. We focus on scientific question answering because scientific answers are high-stakes, evidence-dependent, and naturally compatible with citation-grounded evaluation. We use QASPER \citep{dasigi2021qasper}, which provides research-paper questions, answers, and gold evidence spans, allowing us to evaluate both retrieval and final answers.

As a concrete example, the system may receive a paper and a question such as ``Do they report results only on English data?'' The desired behavior is not just to produce a fluent answer, but to retrieve the right evidence passages and answer in a way that is correct and grounded in the paper. Technically, our task is: given a paper $p$ and a question $q$, retrieve a ranked list of evidence chunks $c_{1:k}$ from $p$ and generate an answer $\hat{y}$, optionally with citations back to those chunks.

This problem is interesting because a successful system would make local question-answering over scientific documents much more useful in settings where users need evidence-backed answers, not just plausible text. The problem is also hard, and simple solutions fail in predictable ways. Lexical BM25 retrieval misses paraphrases and synonymy, while dense retrieval can return topically related but non-evidential text. On the generation side, our own parametric-only baseline performs much worse than retrieval-based systems, and citation forcing alone raises citation frequency without guaranteeing grounded attribution. These failure modes are exactly why our project emphasizes controlled retrieval baselines and grounding-aware evaluation rather than fluent generation alone.

This project draws most directly on course material from the transformer and LLM lectures (attention, pretraining, prompting, and local adaptation), the question answering and information retrieval lecture, and the evaluation lectures on metrics and best practices. Several engineering components in our system were not covered directly in lecture, including FAISS indexing, reciprocal-rank fusion, cross-encoder reranking, and citation-grounded RAG evaluation. We therefore learned those pieces from project readings, model documentation, and system-building practice. We explicitly note these gaps because they shape the remaining work before the final report.

\section{Description of Progress}

\subsection{Methods}

We implemented a fully local RAG pipeline around QASPER. The data layer normalizes raw papers, questions, answers, and evidence into a shared internal schema. We then build deterministic processed splits in a shared SCC course directory so the entire team uses identical artifacts.

For chunking, we implemented the three strategies proposed in our project plan: fixed-size chunks (256 tokens, 32 overlap), sliding-window chunks (512 tokens, 128 overlap), and section-aware chunks aligned to paper structure. Our retrieval stack combines four standard IR components. BM25 \citep{robertson2009bm25} is a sparse lexical ranker,
\[
\resizebox{\columnwidth}{!}{$
\mathrm{BM25}(q,d)=\sum_{t\in q}\mathrm{IDF}(t)\frac{f(t,d)(k_1+1)}{f(t,d)+k_1\!\left(1-b+b\frac{|d|}{\overline{\mathrm{dl}}}\right)},
$}
\]
which is simple and robust but can miss paraphrases and synonymy. Dense retrieval \citep{karpukhin2020dpr,xiao2023bge} embeds the query and chunk and scores them by cosine similarity,
\[
s_{\mathrm{dense}}(q,d)=\cos(e_q,e_d)=\frac{e_q^\top e_d}{\lVert e_q\rVert \lVert e_d\rVert},
\]
which better captures semantics but is more expensive and can retrieve topically similar yet non-evidential text. Hybrid retrieval fuses sparse and dense rankings with reciprocal-rank fusion \citep{cormack2009rrf},
\[
\mathrm{RRF}(d)=\sum_{m\in M}\frac{1}{k+\mathrm{rank}_m(d)},
\]
giving a practical balance between lexical and semantic matching. Our final retrieval stage is a cross-encoder reranker \citep{nogueira2019bert}, which jointly scores each query--chunk pair; it often improves the top few passages, but it is the slowest stage.

For generation, we use the instruction-tuned checkpoint \texttt{microsoft/Phi-3.5-mini-instruct} \citep{microsoft2024phi3}, not a base Phi model. That design choice matters because the task requires explicit answer-format instructions such as yes/no/unanswerable normalization and optional citation formatting. We compare two prompt styles: baseline answering and citation forcing, where answers must cite retrieved passages. Several evaluation metrics were not covered directly in class, so we define them explicitly. If $c$ is the number of overlapping normalized answer tokens, token-level precision and recall are $P=c/|\hat{y}|$ and $R=c/|y|$, and token F1 is $2PR/(P+R)$. BERTScore uses contextual token alignments and reports the same precision/recall/F1 form over embedding matches rather than exact tokens. Retrieval Recall@5 is the fraction of gold evidence spans recovered in the top five chunks, $\mathrm{MRR}@10=1/r$ when the first matching evidence chunk appears at rank $r\le 10$ and $0$ otherwise, and evidence hit is a binary indicator that any gold evidence appears in the top five. Citation precision is $|C_{\mathrm{sup}}|/|C|$, where $C$ is the set of cited chunks and $C_{\mathrm{sup}}$ are the cited chunks matching any gold evidence span. NLI citation precision is $|S_{\mathrm{ent}}|/|S|$, where $S$ is the set of cited claim sentences and $S_{\mathrm{ent}}$ are the cited sentences entailed by at least one cited chunk. Supported-answer rate is the mean of the binary variable \texttt{answer\_supported}, and our hallucination proxy is the mean of $\mathbf{1}[\text{token F1}=0 \lor \texttt{answer\_supported}=0]$. For gold \texttt{unanswerable} questions, we explicitly mark an \texttt{unanswerable} prediction as supported.

\subsection{Experiments}

We completed the retrieval experiments on validation section chunks and then ran the proposal's core ablation conditions A--E as controlled, one-factor-at-a-time comparisons. Condition A starts from fixed chunking, dense retrieval, and baseline prompting. Condition B keeps fixed chunking and the same baseline prompt, but swaps dense retrieval for hybrid retrieval, isolating the effect of retrieval fusion. Condition C keeps hybrid retrieval and baseline prompting but replaces fixed chunks with section-aware chunks, isolating the value of paper structure. Condition D adds the cross-encoder reranker on top of C to test whether better top-of-list ordering improves generation. Condition E keeps D's retrieval pipeline and changes only the prompt to citation forcing, isolating the cost and benefit of explicit attribution. We also added the proposal baselines: parametric-only generation (no retrieval), BM25-only retrieval, and random retrieval over five passages. All generation conditions above were first run on a 100-question validation sample to keep experiments comparable on SCC. After stabilizing the pipeline, we extended A--E, parametric-only, and BM25-only to the full 1005-question validation split and added an NLI-based citation scorer for citation-forcing runs. In addition, we ran exploratory model comparisons with Qwen2.5-7B and Qwen3-4B, but those are side experiments rather than the main midway result story.

\subsection{Results}

Table~\ref{tab:retrieval} summarizes the retrieval-only experiments. Hybrid retrieval gives the best overall validation balance, while hybrid plus reranking improves Recall@5 slightly but does not dominate every metric.

\begin{table}[H]
\centering
\small
\resizebox{\columnwidth}{!}{%
\begin{tabular}{lccc}
\toprule
Method & Recall@5 & MRR@10 & Evidence hit \\
\midrule
BM25 & 0.5769 & 0.4110 & 0.7152 \\
Dense & 0.5886 & 0.4820 & 0.7238 \\
Hybrid & 0.6313 & 0.5097 & 0.7778 \\
Hybrid + rerank & 0.6410 & 0.5027 & 0.7702 \\
\bottomrule
\end{tabular}
}
\caption{Validation retrieval results on section-aware chunks.}
\label{tab:retrieval}
\end{table}

Table~\ref{tab:core} shows the main A--E generation results. The strongest condition so far is C: section-aware chunking with hybrid retrieval and baseline prompting. Moving from fixed hybrid retrieval (B) to section-aware hybrid retrieval (C) gives the largest gain in the main table. Reranking (D) does not beat C on this 100-question sample, although it does slightly improve supported-answer rate. Citation forcing (E) produces citations very reliably, but still lowers answer quality and grounding relative to the strongest non-citation setting.

\begin{table}[H]
\centering
\small
\resizebox{\columnwidth}{!}{%
\begin{tabular}{lcccccc}
\toprule
Cond. & Retrieval & Chunking & Prompt & Token F1 & BERTScore F1 & Halluc. \\
\midrule
A & Dense & Fixed & Baseline & 0.3636 & 0.8918 & 0.50 \\
B & Hybrid & Fixed & Baseline & 0.3365 & 0.8834 & 0.52 \\
C & Hybrid & Section & Baseline & \textbf{0.3947} & \textbf{0.8946} & \textbf{0.39} \\
D & Hybrid + rerank & Section & Baseline & 0.3697 & 0.8910 & 0.43 \\
E & Hybrid + rerank & Section & Citation & 0.3569 & 0.8871 & 0.54 \\
\bottomrule
\end{tabular}
}
\caption{Core proposal ablations on a 100-question validation sample. Condition E is the citation-forcing setting, while Condition C is the strongest non-citation condition in this table.}
\label{tab:core}
\end{table}

Table~\ref{tab:baselines} gives the three extra baselines from the proposal. Parametric-only generation is much worse than any retrieval-based condition, which supports the central claim that retrieval materially reduces hallucination in this task. BM25-only performs better than random retrieval and nearly matches the weaker dense/hybrid ablations on this sample, while random retrieval remains clearly worse on grounding.

\begin{table}[H]
\centering
\small
\resizebox{\columnwidth}{!}{%
\begin{tabular}{lcccc}
\toprule
Baseline & Token F1 & BERTScore & Halluc. & Supported \\
\midrule
Parametric-only & 0.1600 & 0.8397 & 0.95 & 0.05 \\
BM25-only & 0.3426 & 0.8896 & 0.49 & 0.70 \\
Random retrieval & 0.3257 & 0.8856 & 0.68 & 0.46 \\
\bottomrule
\end{tabular}
}
\caption{Additional baselines on the same 100-question validation sample.}
\label{tab:baselines}
\end{table}

One practical lesson from these experiments is that automatic similarity metrics need to be interpreted carefully. BERTScore is useful as a secondary paraphrase signal, but it can overstate answer quality when two texts sound related without matching the specific gold answer. For that reason, we treat token F1, supported-answer rate, hallucination rate, and later manual judgments as stronger evidence for the main conclusions than BERTScore alone.

Table~\ref{tab:nli-followup} summarizes the later full-validation follow-up on the strongest Phi settings. The best full-validation baseline is now D rather than C, but the overall story is unchanged: retrieval plus reranking helps the baseline system, while citation forcing still trades answer support for explicit attribution. The NLI-based citation result is especially important because it shows that many cited claims are still not actually entailed by the cited evidence. This confirms the error-analysis finding that explicit citation formatting alone is not enough to ensure grounded attribution.

\begin{table}[H]
\centering
\small
\resizebox{\columnwidth}{!}{%
\begin{tabular}{lccccc}
\toprule
Setting & Questions & Token F1 & Cite Prec. & Cite Rate & NLI Cite \\
\midrule
D full val. & 1005 & \textbf{0.3238} & -- & -- & -- \\
E full val. & 1005 & 0.3215 & 0.3230 & 0.9274 & 0.1416 \\
\bottomrule
\end{tabular}
}
\caption{Full-validation follow-up on the strongest Phi settings. Citation-only columns are not applicable to D because baseline prompting does not produce citations.}
\label{tab:nli-followup}
\end{table}

\section{Project Plan}

Before the final report, we plan to do three things. First, we will consolidate the now-complete full-validation A--E table, the supporting baselines, and our completed small human-evaluation pass into the final writeup. Second, we will decide whether one additional verifier-gated citation experiment is justified now that the NLI citation check is available. Third, we will update the final report to clearly separate our core Phi-based ablation story from the exploratory Qwen side experiments.

We also want to be explicit about what is still outside the core proposal scope. We have not yet run the proposal's stretch-generator conditions with Llama-3-8B or Mistral-7B, and we have not yet added a second local faithfulness judge beyond the current NLI-based attribution scorer. These missing pieces are not due to a conceptual shift in the project; they are deferred because the main effort went into building a stable local SCC pipeline, completing the full-validation Phi ablations, and adding grounded-attribution evaluation.

\section{Conclusion}

By the midway point, the core system is working end-to-end and the main proposal ablation table is substantially complete. The clearest result so far is that chunking and retrieval choices matter a great deal: section-aware chunking with hybrid retrieval gives the best initial sample results, while the later full-validation follow-up favors hybrid retrieval with reranking under baseline prompting. Citation forcing is promising for attribution behavior, but the new NLI result shows that frequent citation formatting still does not guarantee grounded attribution. The remaining work is therefore focused less on infrastructure and more on tightening evaluation and improving grounded generation quality.

{\small
\bibliographystyle{acl_natbib}
\bibliography{references}
}

\end{document}
